% !TeX root = ../../../main.tex
\FloatBarrier
\subsection{Dataset-stratified performance}\label{ch08:supp:dataset-stratified}
\noindent
\textbf{Methods.} We repeated the primary evaluation separately for the 15 i2b2-derived cases and 25 MIMIC-IV cases using event-matching threshold 0.1. Event match rate, concordance, and AULTC were recomputed within each source and reported separately for unimodal and multimodal reconstruction. Source-specific 95\% confidence intervals were obtained by resampling cases with replacement within each source using the same 200-resample case-level bootstrap procedure as the marginal confidence intervals in the main analysis. These tables provide descriptive source-specific estimates; they are not paired tests of the multimodal-minus-unimodal difference.

\noindent
\textbf{Results.} Multimodal calibration changed event match rate little within either source. Performance was generally higher on the i2b2-derived subset: five of six models had higher event match rates and all six had higher AULTC than on MIMIC-IV. Temporal shifts were heterogeneous. On i2b2, multimodal calibration produced mostly small AULTC gains and modest concordance changes. On MIMIC-IV, multimodal temporal point estimates were generally unchanged or lower in the supplied source-stratified outputs, with the largest decreases for DeepSeek V3.2 and Qwen3.5-397B-A17B. The broad overlap of the marginal intervals warrants cautious interpretation.

\par\medskip
\noindent
% !TeX root = ../../../main.tex
\begin{table}[htbp]
\caption{Performance on the i2b2-derived subset (n = 15) at event-matching threshold 0.1 for unimodal and multimodal timeline reconstruction.}\ChEightFloatLabel{ch08:tab:supp-i2b2}
\ChEightTableBodyStart
\begin{ChEightFloatBody}


\centering
\scriptsize
\begin{adjustbox}{max width=\textwidth}
\begin{tabular}{@{}lcccccc@{}}
\toprule
& \multicolumn{2}{c}{\textbf{Event match rate}} & \multicolumn{2}{c}{\textbf{Concordance}} & \multicolumn{2}{c}{\textbf{AULTC}} \\
\cmidrule(lr){2-3}\cmidrule(lr){4-5}\cmidrule(l){6-7}
\textbf{Model} & \textbf{Unimodal} & \textbf{Multimodal} & \textbf{Unimodal} & \textbf{Multimodal} & \textbf{Unimodal} & \textbf{Multimodal} \\
\midrule
DeepSeek R1 & 0.51 [0.42, 0.61] & 0.51 [0.42, 0.61] & 0.77 [0.71, 0.84] & 0.77 [0.70, 0.83] & 0.82 [0.80, 0.85] & 0.82 [0.80, 0.85] \\
KimiK2-Instruct & 0.40 [0.33, 0.50] & 0.40 [0.32, 0.50] & 0.73 [0.69, 0.81] & 0.74 [0.70, 0.80] & 0.78 [0.75, 0.80] & 0.78 [0.75, 0.81] \\
GPT-OSS-120B & 0.53 [0.48, 0.60] & 0.53 [0.48, 0.60] & \textbf{0.80 [0.76, 0.84]} & 0.77 [0.75, 0.83] & 0.80 [0.77, 0.83] & 0.81 [0.79, 0.84] \\
DeepSeek V3.2 & \textbf{0.62 [0.55, 0.70]} & \textbf{0.62 [0.55, 0.70]} & 0.77 [0.73, 0.82] & 0.79 [0.70, 0.82] & 0.82 [0.79, 0.84] & 0.82 [0.79, 0.85] \\
GLM5 & 0.28 [0.21, 0.38] & 0.28 [0.21, 0.38] & 0.80 [0.72, 0.88] & \textbf{0.81 [0.76, 0.85]} & \textbf{0.83 [0.79, 0.86]} & \textbf{0.84 [0.80, 0.87]} \\
Qwen3.5-397B-A17B & 0.61 [0.54, 0.68] & 0.61 [0.54, 0.68] & 0.78 [0.75, 0.84] & 0.78 [0.72, 0.85] & 0.80 [0.78, 0.83] & 0.81 [0.79, 0.84] \\
\bottomrule
\end{tabular}
\end{adjustbox}
\par
\begin{minipage}{\textwidth}
\footnotesize\textit{Note:} Values are point estimates with source-specific case-level bootstrap 95\% confidence intervals in brackets. Bold indicates the largest unrounded point estimate within each metric-modality column. All displayed values are rounded to two decimals.
\end{minipage}

\end{ChEightFloatBody}
\end{table}
\par\medskip

\par\medskip
\noindent
% !TeX root = ../../../main.tex
\begin{table}[htbp]
\caption{Performance on the MIMIC-IV subset (n = 25) at event-matching threshold 0.1 for unimodal and multimodal timeline reconstruction.}\ChEightFloatLabel{ch08:tab:supp-mimiciv}
\ChEightTableBodyStart
\begin{ChEightFloatBody}


\centering
\scriptsize
\begin{adjustbox}{max width=\textwidth}
\begin{tabular}{@{}lcccccc@{}}
\toprule
& \multicolumn{2}{c}{\textbf{Event match rate}} & \multicolumn{2}{c}{\textbf{Concordance}} & \multicolumn{2}{c}{\textbf{AULTC}} \\
\cmidrule(lr){2-3}\cmidrule(lr){4-5}\cmidrule(l){6-7}
\textbf{Model} & \textbf{Unimodal} & \textbf{Multimodal} & \textbf{Unimodal} & \textbf{Multimodal} & \textbf{Unimodal} & \textbf{Multimodal} \\
\midrule
DeepSeek R1 & 0.36 [0.28, 0.45] & 0.36 [0.28, 0.45] & 0.70 [0.68, 0.80] & 0.70 [0.66, 0.80] & \textbf{0.79 [0.75, 0.82]} & \textbf{0.78 [0.75, 0.82]} \\
KimiK2-Instruct & 0.22 [0.19, 0.26] & 0.21 [0.19, 0.25] & 0.71 [0.59, 0.82] & 0.67 [0.56, 0.82] & 0.64 [0.56, 0.71] & 0.63 [0.56, 0.70] \\
GPT-OSS-120B & 0.27 [0.24, 0.34] & 0.27 [0.23, 0.34] & 0.70 [0.64, 0.75] & 0.70 [0.65, 0.77] & 0.76 [0.72, 0.80] & 0.74 [0.69, 0.79] \\
DeepSeek V3.2 & \textbf{0.46 [0.38, 0.55]} & \textbf{0.46 [0.38, 0.55]} & 0.77 [0.70, 0.83] & 0.71 [0.62, 0.77] & 0.76 [0.71, 0.81] & 0.67 [0.58, 0.76] \\
GLM5 & 0.34 [0.26, 0.41] & 0.33 [0.26, 0.40] & \textbf{0.82 [0.77, 0.85]} & \textbf{0.80 [0.73, 0.83]} & 0.77 [0.71, 0.83] & 0.76 [0.71, 0.82] \\
Qwen3.5-397B-A17B & 0.35 [0.24, 0.48] & 0.32 [0.22, 0.45] & 0.81 [0.78, 0.83] & 0.78 [0.67, 0.81] & 0.75 [0.72, 0.77] & 0.70 [0.65, 0.75] \\
\bottomrule
\end{tabular}
\end{adjustbox}
\par
\begin{minipage}{\textwidth}
\footnotesize\textit{Note:} Values are point estimates with source-specific case-level bootstrap 95\% confidence intervals in brackets. Bold indicates the largest unrounded point estimate within each metric-modality column. All displayed values are rounded to two decimals.
\end{minipage}

\end{ChEightFloatBody}
\end{table}
\par\medskip

